% !TEX program = xelatex
% !TEX options = -synctex=1 -interaction=nonstopmode -file-line-error
% example-manuscript.tex
% Demonstrates rossmanuscript: title page, structured abstract,
% sections with statistical results, and disclosure statements.

\begin{filecontents}{example-manuscript.bib}
@article{smith2024sleep,
  author  = {Smith, J. A. and Lee, K. R.},
  title   = {Sleep architecture and cognitive decline in older adults},
  journal = {Sleep},
  year    = {2024},
  volume  = {47},
  number  = {3},
  pages   = {1--12},
  doi     = {10.1093/sleep/zsae012},
}
@article{jones2023mediators,
  author  = {Jones, P. B. and Patel, M.},
  title   = {Affective mediators of the sleep--cognition relationship},
  journal = {Journal of Clinical Psychology},
  year    = {2023},
  volume  = {79},
  pages   = {455--470},
  doi     = {10.1002/jclp.23498},
}
\end{filecontents}

\documentclass[palette=lancet, journal={Sleep}]{rossmanuscript}
\addbibresource{example-manuscript.bib}

% ── Metadata ─────────────────────────────────────────────────────
\title{Sleep Architecture, Affective Symptoms, and Domain-Specific
       Cognition: A Mediation Analysis}
\shorttitle{Sleep, Affect, and Cognition}
\author{Ross A.\ Dunne\textsuperscript{1},
        Sarah K.\ Mitchell\textsuperscript{2}}
\affiliation{%
  \textsuperscript{1}Department of Psychology, University of Limerick\\
  \textsuperscript{2}School of Medicine, Trinity College Dublin}
\orcid{0000-0002-1234-5678}
\correspondingauthor{Ross A.\ Dunne}{ross.dunne@ul.ie}
\authornote{This study was pre-registered on OSF (osf.io/abcde).}
\keywords{sleep architecture, cognition, mediation, affect, ageing}

% ── Disclosures ──────────────────────────────────────────────────
\conflictofinterest{The authors declare no conflicts of interest.}
\funding{This work was supported by the Irish Research Council (grant GOIPG/2024/1234).}
\dataavailability{De-identified data are available at \url{https://osf.io/abcde}.}

\begin{document}
\maketitlepage

\structuredabstract
  {Poor sleep is associated with cognitive decline, yet the mediating role
   of affective symptoms remains unclear \autocite{smith2024sleep}.}
  {Community-dwelling older adults ($N = 312$, age 60--85) completed
   polysomnography, the DASS-21, and the RBANS. Parallel mediation models
   tested indirect effects of sleep architecture on five cognitive domains
   through depressive and anxiety symptoms.}
  {Slow-wave sleep percentage predicted delayed memory via depressive
   symptoms (\betacoef{-0.18}; \ci{-0.28}{-0.34}{-0.08}; \pval{0.002}).
   No significant indirect effect emerged through anxiety
   (\pval{0.41}).}
  {Depressive symptoms partially mediate the relationship between
   reduced slow-wave sleep and memory decline, suggesting a potential
   therapeutic target \autocite{jones2023mediators}.}

\section{Introduction}

Sleep disturbance is increasingly recognised as a modifiable risk factor
for cognitive decline in older adults. Polysomnographic studies have linked
reductions in slow-wave sleep (SWS) to poorer memory consolidation, yet the
mechanisms underlying this relationship are not fully understood.

\section{Method}

\subsection{Participants}

We recruited 312 community-dwelling adults aged 60--85 years ($M = 71.4$,
$SD = 6.2$; 58\% female) from primary-care registries. Exclusion criteria
included diagnosed dementia, untreated sleep apnoea, and current
psychotropic medication use.

\subsection{Measures}

Cognitive function was assessed with the Repeatable Battery for the
Assessment of Neuropsychological Status (RBANS). Sleep architecture was
derived from single-night in-lab polysomnography.

\section{Results}

A significant total effect of SWS\% on delayed memory was observed
(\tstat{308}{-3.14}; \pval{0.002}). The indirect effect through depressive
symptoms was significant (\betacoef{-0.18}; \CI{95}{-0.18}{-0.34}{-0.08}).
The model explained 22\% of the variance in delayed memory
(\fstat{3}{308}{28.97}; \pval{0.00001}).

\section{Discussion}

These findings extend prior work by identifying depressive symptoms as a
partial mediator of the sleep--memory relationship. Clinically, targeting
mood disturbance in older adults with sleep complaints may yield
downstream cognitive benefits.

\printdisclosures

\printbibliography

\end{document}
